\documentclass[12pt]{article}
\usepackage[utf8]{inputenc}
\usepackage[T1]{fontenc}
\usepackage{hyperref}
\title{Tolowa Language Resources for K-12 Curriculum Development}
\author{Compiled by Claude Sonnet 5 (claude-sonnet-5) for Andr\'as Kornai}
\date{}
\begin{document}
\maketitle

\section{Introduction}
Tolowa (Tolowa Dee-ni') is a dormant Athabaskan language of the Del Norte County area in northern California and southern Oregon. The Tolowa Dee-ni' Nation runs revitalization work through its Wee-ya' (Language) Program. Search coverage was limited, and only a few resources were seen.

\section{Primary Educational Resources}
\begin{itemize}
\item Tolowa Dee-ni' Nation Wee-ya' Program: its stated aim is to restore Tolowa as a spoken language. Activities include community and high school classes, a language-in-the-home initiative, apps and curriculum development: \url{https://www.weeyadvn.com/about/}
\item \emph{Now You're Speaking Tolowa} (catalogue entry): \url{https://wiki.ubc.ca/Documentation:RelLex/Now_You%27re_Speaking_Tolowa}
\end{itemize}

\section{Linguistic and Reference Materials}
\begin{itemize}
\item Tolowa Dee-ni' Nation online dictionary (catalogue entry): \url{https://wiki.ubc.ca/Documentation:RelLex/Tolowa_Dee-ni%27_Wee-ya%27-dvn}
\item Wee-ya' grammar workshop slides (2018 Grammar Day 3): \url{https://www.weeyadvn.com/userfiles/2018-Grammar-Day-3-powerpoint.pdf}
\item Reference article: \url{https://en.wikipedia.org/wiki/Tolowa_language}
\item NSF-funded lexicon project with sound files, a collaboration with Smith River Rancheria, the University of Oregon and the Del Norte County Unified School District: \url{https://grantome.com/grant/NSF/BCS-0854556}
\end{itemize}

\section{Religious and Community Resources}
The Wee-ya' Committee of seven tribal citizens advises the program, per the Nation's site above. A documentary on the revival, \emph{Still, We Live On}: \url{https://now.humboldt.edu/news/still-we-live-film-showcases-tolowa-language-resilience-and-revival}

\section{Digital and Audio Resources}
The program reportedly offers iOS and Android apps and an online dictionary. Audio exists in the NSF lexicon project, but I did not verify where it is hosted.

\section{Curriculum Development Context}
Sources disagree on the language's status. Some describe it as dormant, others as having a few remaining speakers, and all describe active revival. Any curriculum should be developed with the Tolowa Dee-ni' Nation's Wee-ya' Program.
\end{document}